\documentclass[10pt,a4paper]{article}
\usepackage[margin=0.75in,top=0.65in,bottom=0.65in,headheight=14pt]{geometry}
\usepackage[T1]{fontenc}
\usepackage{lmodern,microtype,graphicx,booktabs,array,tabularx,amsmath,amssymb}
\usepackage{caption,float,placeins,xcolor,setspace,enumitem,fancyhdr,needspace}
\usepackage{hyperref,titlesec}
\definecolor{ink}{HTML}{000000}
\definecolor{accent}{HTML}{153E6F}
\definecolor{rust}{HTML}{BB603A}
\hypersetup{colorlinks=true,linkcolor=accent,citecolor=accent,urlcolor=accent,
 pdftitle={Food Inflation and Unemployment in India},pdfauthor={K Harshavardhan and Pasula Kohith}}
\urlstyle{same}
\graphicspath{{assets/}}
\setstretch{1.16}
\setlength{\parindent}{0pt}
\setlength{\parskip}{5pt}
\setlength{\tabcolsep}{5pt}
\setlength{\intextsep}{8pt plus 2pt minus 2pt}
\setlength{\textfloatsep}{8pt plus 2pt minus 2pt}
\setlength{\emergencystretch}{1em}
\renewcommand{\arraystretch}{1.16}
\captionsetup{font=small,labelfont=bf,skip=5pt}
\titleformat{\section}{\Large\bfseries\color{ink}}{\thesection}{0.6em}{}
\titleformat{\subsection}{\large\bfseries\color{ink}}{\thesubsection}{0.6em}{}
\titlespacing*{\section}{0pt}{12pt}{6pt}
\titlespacing*{\subsection}{0pt}{9pt}{4pt}
\clubpenalty=10000\widowpenalty=10000\displaywidowpenalty=10000
\raggedbottom
\setlist{leftmargin=1.3em,itemsep=2pt,topsep=3pt}
\pagestyle{fancy}\fancyhf{}
\fancyhead[L]{\small Food inflation and unemployment in India}
\fancyhead[R]{\small HS455}
\fancyfoot[C]{\small\thepage}
\renewcommand{\headrulewidth}{0.3pt}
\newcommand{\tabnote}[1]{\par\vspace{2pt}{\footnotesize\emph{Note:} #1\par}}
\newcommand{\source}[1]{{\footnotesize\emph{Source:} #1\par}}
\newcommand{\fig}[4]{\begin{figure}[H]\centering\includegraphics[width=\linewidth]{#1}\caption{#2}\label{#3}\source{#4}\end{figure}}
\begin{document}
\hypersetup{pageanchor=false}
\begin{titlepage}
\thispagestyle{empty}
\begin{center}
\vspace*{0.3cm}
{\fontsize{20}{26}\selectfont\bfseries Food Inflation and Unemployment:\\[6pt] Evidence from Urban India}\\[0.6cm]
{\Large\bfseries Headline Inflation, Core Inflation and Food-Price Variability}\\[0.8cm]
{\Large\itshape A Project Report Submitted}\\[0.25cm]
{\Large\itshape for the Course}\\[0.4cm]
{\LARGE\bfseries HS455}\\[0.3cm]
{\large\itshape Instructors}\\[0.15cm]
{\large\bfseries Dr. Smruti Ranjan Behera}\\[0.1cm]
{\large\bfseries Samaresh Bardhan\qquad Bhavesh Garg}
\vfill
\includegraphics[width=0.44\textwidth]{iit_ropar_logo.png}\\
\vfill
{\Large\bfseries HS455 COURSE PROJECT}\\[0.25cm]
{\Large\bfseries INDIAN INSTITUTE OF TECHNOLOGY ROPAR}\\[0.25cm]
{\large\bfseries November 2026}\\[0.25cm]
{\large\itshape by}\\[0.25cm]
{\Large\bfseries Students}\\[0.25cm]
{\renewcommand{\arraystretch}{1.3}
\begin{tabular}{p{6.1cm}p{4.6cm}}\toprule
\textbf{Name} & \textbf{Entry No.}\\\midrule
K Harshavardhan & 2023MCB1216\\
Pasula Kohith & 2023EEB1232\\\bottomrule
\end{tabular}}
\end{center}
\end{titlepage}
\setcounter{page}{1}\hypersetup{pageanchor=true}
% MAIN_TEXT_START
\section{Background and motivation}

This report studies the relationship between inflation and unemployment in India and measures food's contribution to headline-inflation variability. Demand growth can increase prices and employment together. A food-supply disruption can instead raise prices during weak employment conditions. Separating food from the remaining CPI groups helps us examine whether these different sources of price change affect the estimated unemployment relationship.

The Phillips curve suggests a negative short-run relationship: lower unemployment may accompany stronger wage and price pressures. In actual data, however, demand, supply shocks, expectations and policy responses operate together. Excluding food might reveal a clearer negative relationship if food-price fluctuations obscure demand-related inflation. We test this possibility using the same observations for headline and core inflation.

India is a useful case because Food and Beverages carry 45.86\% of the 2012-base Combined Consumer Price Index (CPI), compared with 36.29\% in Urban CPI. These expenditure weights are large, but they do not directly measure food's contribution to inflation variability. That contribution also depends on price movements and their covariance with other components.

Three studies inform the analysis. Behera, Wahi and Kapur find a role for demand conditions in state-level inflation, with a stronger response for core inflation in their output-gap framework~\cite{behera2017}. Their activity measure differs from our unemployment rate. Patra, John and George investigate food--non-food dynamics using an error-correction model~\cite{patra2024}, while Mohan and co-authors examine inflation volatility and vegetable prices~\cite{mohan2026}. Our estimates use simpler models and a different sample, so they are compared with these studies rather than treated as replications.

A distinction between the price level and inflation is useful here. A household can still face expensive food after food inflation falls, because a slower rate of increase does not reverse earlier price rises. Likewise, a low average inflation rate can coexist with sharp monthly fluctuations. Our unemployment exercise concerns the rate of price change, while the variance exercise concerns fluctuations around its average. Neither calculation measures the complete burden on a household budget.

We estimate the unemployment relationship separately from the variance allocation. A large food share in headline variability does not, by itself, establish that excluding food improves the Phillips-curve relationship. Each question therefore has its own measure, sample and hypothesis. The conclusions depend on those calculations rather than on whether the estimated signs match the expected pattern.

Our approach joins an Urban unemployment comparison with a broader CPI accounting exercise. The first asks whether core inflation reveals the expected negative association. The second measures how much headline variability is allocated to food. A monthly predictability exercise then extends the analysis to food--non-food spillovers. The three exercises address different questions and should not be read as interchangeable evidence.

\section{Research questions and hypotheses}

\textbf{Question 1: Does excluding food strengthen the negative inflation--unemployment relationship?} We compare headline and operational-core inflation on identical quarterly observations. If their unemployment slopes are $\beta_H$ and $\beta_C$, the hypothesis is a more negative core slope, $\delta=\beta_C-\beta_H<0$, together with a negative core association. Correlations provide a second descriptive comparison. A stronger positive correlation would not support the proposed negative relationship.

\textbf{Question 2: How much variation in headline inflation is due to food?} We expect food's variance allocation to exceed its expenditure weight. ``Due to'' is interpreted as an accounting allocation, including covariance with other components, rather than the share caused by independently identified food shocks.

\textbf{Extension: Does food inflation predict subsequent non-food or core inflation?} We examine whether three delayed food-inflation coefficients are jointly informative after accounting for persistence and calendar effects. This tests conditional predictability; it cannot establish transmission through wages, costs or expectations.

\clearpage
\section{Data and methodology}
\subsection{Sources, sample and comparability}

The price data are monthly 2012-base CPI levels, January 2013--December 2025~\cite{cpiapi}. The unemployment series is the Periodic Labour Force Survey (PLFS), All-India Urban, persons aged 15 and above, Current Weekly Status (CWS). CWS uses the preceding seven days; the unemployment rate measures unemployed persons as a percentage of the labour force~\cite{plfs2024}. There are 27 matched quarters from 2018Q2 to 2024Q4.

We match Urban prices with Urban unemployment. Combined CPI, covering rural and urban households, is used separately for the wider inflation-composition question. This avoids silently treating Urban labour-market evidence as national evidence. The redesigned PLFS begins in January 2025~\cite{plfschange}; it is not joined to the earlier quarterly series. Later-base CPI observations are also kept separate.

Data checks found two unemployment API discrepancies. For 2019Q2, 21.6\% refers to the youth row rather than the required 15+ rate of 8.8\%. For 2019Q3, the API gives 8.1\%, while the published bulletin gives 8.3\%~\cite{plfs2019a,plfs2019b}. We retain the bulletin values and record both corrections. The other 25 observations agree with the bulletin transcription.

\subsection{Price-basket construction and inflation measures}

Headline uses the official General index. Operational core excludes Food and Beverages and Fuel and Light, then rescales the retained group weights to sum to one. Food-only exclusion is also tested, since the assignment specifically asks about removing food.

\begin{table}[H]\centering\small
\caption{Urban CPI weights and the operational-core basket}\label{tab:weights}
\begin{tabularx}{\linewidth}{Xrrl}\toprule
Group & Urban weight (\%) & Share of core (\%) & Core treatment\\\midrule
Food and beverages &36.29&--&Exclude\\
Pan, tobacco and intoxicants &1.36&2.34&Retain\\
Clothing and footwear &5.57&9.58&Retain\\
Housing &21.67&37.28&Retain\\
Fuel and light &5.58&--&Exclude\\
Miscellaneous &29.53&50.80&Retain\\\bottomrule
\end{tabularx}
\tabnote{Retained weights sum to 58.13\%. Vehicle fuels and precious metals remain within miscellaneous items.}\source{MoSPI CPI release tables~\cite{cpirelease}; calculations.}
\end{table}

For December 2024, the retained indices give a worked example:
\begin{equation}
C_{Dec24}=\frac{1.36(212.2)+5.57(187.8)+21.67(181.7)+29.53(182.0)}{58.13}=183.1505.
\end{equation}
Housing represents 37.28\% of this core, so a construction without housing is checked as well.

Rescaling the retained weights prevents the core index from falling simply because two groups have been removed. It keeps the index on a comparable scale before growth rates are calculated. The construction is nevertheless selective: excluding the Fuel and Light group does not remove every item influenced by energy costs. We therefore report exactly what is retained, and use the food-only comparison to separate the assignment's question from our broader operational definition.

Quarterly levels average all three monthly indices. Year-on-year log inflation is calculated as $100\ln(\bar I_q/\bar I_{q-4})$.
Urban headline means of 183.7333 in 2023Q4 and 192.9667 in 2024Q4 imply inflation of $4.9032\%$, calculated as $100\ln(192.9667/183.7333)$. This is not an average of rounded published monthly inflation rates.

\clearpage
\subsection{Descriptive statistics and regression estimation}

We begin with means, medians, sample standard deviations, ranges, skewness and excess kurtosis. The mean describes the average quarter; the median is less affected by an extreme observation. Standard deviation measures dispersion. Positive skewness indicates a longer upper tail; excess kurtosis describes tail weight relative to a normal distribution, whose population value is zero. We use bias-corrected sample moments and avoid treating a few observations as a reliable distributional test.

The standard deviation uses the squared distance of each observation from the sample mean, divided by one less than the number of observations, before taking the square root. Squaring gives a large observation more influence than the median does. Skewness and kurtosis go further by using third and fourth powers of deviations. They are useful warning signals about unusual observations, but can themselves be unstable in a short sample. We therefore read them with the actual plotted observations rather than classify the data from a single cutoff. An extreme value is checked against its source before being treated as either a genuine event or an error; statistical influence alone is not grounds for correcting it.

Pearson correlation measures linear association. Ordinary least squares (OLS) then estimates an intercept and unemployment slope for each inflation measure:
\begin{equation}
\pi_{H,q}=\alpha_H+\beta_H U_q+e_{H,q},\qquad
\pi_{C,q}=\alpha_C+\beta_C U_q+e_{C,q}.
\label{eq:reg}
\end{equation}
The slope is an inflation percentage-point difference associated with a one-percentage-point difference in unemployment. Regressing $D_q=\pi_{C,q}-\pi_{H,q}$ on the same intercept and $U_q$ estimates $\delta$ directly. This uses the relationship between the two inflation outcomes, rather than comparing their separate significance levels.

A second transformation is seasonally adjusted quarter-on-quarter annualised log inflation (q/q SAAR). Robust STL removes a seasonal component from monthly log indices, using period 12 and seasonal window 13. Adjusted logs are converted back to levels before quarterly averaging; inflation is $400\ln(\bar I_q^{SA}/\bar I_{q-1}^{SA})$. Annualisation changes the scale, not the forecast horizon. Annual inflation overlaps across quarters, while q/q inflation is more sensitive to seasonal adjustment; both are therefore reported.

The full sample is compared with a fixed exclusion of 2020Q2, 2020Q3 and 2021Q2, leaving 24 quarters. April--May 2020 CPI includes component-specific macro-imputation~\cite{imputation}; the exclusion also captures exceptional pandemic labour-market conditions. We check alternative baskets, a calendar trend, seasonal settings and leave-one-quarter-out estimates.

\subsection{Exact variance accounting and uncertainty}

For this exercise we use ordinary percentage changes so group contributions can add exactly. With official headline index $H_t$, group index $I_{g,t}$, fractional weight $w_g$ and lag $k$,
\begin{equation}
\pi_t=100\left(\frac{H_t}{H_{t-k}}-1\right),\qquad
c_{g,t}=100w_g\frac{I_{g,t}-I_{g,t-k}}{H_{t-k}}.
\end{equation}
We retain a residual $r_t=\pi_t-\sum_g c_{g,t}$ and allocate variance as
\begin{equation}
s_g=100\frac{\operatorname{Cov}(c_g,\pi)}{\operatorname{Var}(\pi)}.
\label{eq:share}
\end{equation}
All component and residual shares sum to 100\%. The covariance allocation allows components to reinforce or offset one another; it is not restricted to positive shares. Twelve-month moving-block resampling gives descriptive uncertainty ranges. Matched-period, imputation-exposure and subperiod comparisons check stability.

Regression uncertainty uses HC3, which allows unequal residual variance and adjusts for leverage, alongside calendar-based HAC checks allowing serial dependence. Short samples and overlapping inflation still weaken calibration. We show nominal ranges without significance stars. Expectations, common supply shocks and policy responses remain outside the simple regression, so it is descriptive rather than causal.

\clearpage
\section{Discussion of results}
\subsection{A first look at the matched Urban series}

Unemployment reaches 20.8\% in 2020Q2 (Figure~\ref{fig:paths}). Neither inflation series shows a stable inverse movement with unemployment over the full period. We therefore examine sample moments, fitted slopes and the influence of individual quarters before interpreting the relationship.

\fig{01_time_paths.pdf}{Matched Urban unemployment and YoY log inflation, 2018Q2--2024Q4. Shading marks 2020--2021.}{fig:paths}{MoSPI/NSO CPI and PLFS; calculations.}

\begin{table}[H]
\centering\small
\caption{Descriptive statistics on matched quarter samples}
\label{tab:descriptive}
\begin{tabular}{lrrrrrrr}
\toprule
Series & $n$ & Mean & Median & SD & Skewness & Excess kurtosis & Min--max\\
\midrule
Unemployment & 27 & 8.88 & 8.30 & 2.98 & 2.69 & 9.45 & 6.40--20.80\\
Headline: YoY & 27 & 5.21 & 5.35 & 1.02 & -0.23 & -0.98 & 3.30--6.78\\
Core: YoY & 27 & 4.84 & 5.29 & 0.99 & -0.44 & -1.12 & 3.02--6.34\\
Headline: q/q SAAR & 27 & 5.20 & 5.10 & 2.14 & 0.03 & 0.51 & 0.05--9.90\\
Core: q/q SAAR & 27 & 4.66 & 4.89 & 1.23 & 0.03 & 0.00 & 2.35--7.41\\
Unemployment: exclusion & 24 & 8.05 & 8.00 & 1.32 & 0.16 & -1.55 & 6.40--10.30\\
\bottomrule
\end{tabular}
\tabnote{All values are percentages or percentage points except skewness and excess kurtosis. SD uses $n-1$; skewness and excess kurtosis use finite-sample bias corrections. The exclusion removes 2020Q2, 2020Q3 and 2021Q2.}
\end{table}


The unemployment mean is 8.88\%, above its median of 8.30\%. Its skewness of 2.69 and excess kurtosis of 9.45 reflect the exceptional upper tail. Excluding the three fixed quarters brings the mean to 8.05\%, the median to 8.00\%, and the standard deviation from 2.98 to 1.32 percentage points. Thus, the pandemic affects both the typical level and the spread available for estimation.

The same arithmetic checks the reported spread: the unemployment squared deviations sum to 231.5667, giving $s_U=\sqrt{231.5667/26}=2.9844$. The median is the fourteenth ordered observation, 8.3. The mean, median and standard deviation describe different features of the same observations.

For annual inflation, headline and core standard deviations are close: 1.02 and 0.99. For q/q SAAR, they are 2.14 and 1.23. Core's short-run standard deviation is about 42.4\% lower, calculated as $100(1-1.2328/2.1388)$. Removing food and fuel smooths short-run inflation; whether it strengthens a negative unemployment relationship remains a separate question.

\clearpage
\subsection{Distribution and correlation evidence}

Figure~\ref{fig:distribution} shows the observed distributions. Most unemployment observations lie far below the pandemic peak. The inflation boxplots show similar annual dispersion, but a wider short-run headline distribution. The negative excess kurtosis of annual headline and core inflation does not establish normality or stability; with only 27 quarters, these moments are descriptive summaries.

\fig{02_distributions.pdf}{Observed distributions. Boxes show quartiles and medians; dots show all observations. H: headline; C: core; q/q: seasonally adjusted annualised inflation.}{fig:distribution}{Matched-quarter datasets; calculations.}

\fig{03_correlations.pdf}{YoY Pearson correlations: full and fixed-exclusion samples. U denotes unemployment. All correlations use matched rows.}{fig:correlation}{Matched Urban CPI and PLFS; calculations.}

The full-sample unemployment correlation is $+0.276$ with headline inflation and $+0.370$ with core. After the exclusion, it becomes $-0.080$ for headline and $+0.537$ for core. A larger core correlation is therefore a closer \emph{positive} association, rather than the stronger negative association proposed in Question 1.

Correlation standardises the covariance by both series' standard deviations. A regression slope instead keeps the original units. In this simple model their link is $\hat\beta=r(s_\pi/s_U)$. Thus a smaller inflation standard deviation can change the slope even when correlation is unchanged. Comparing both summaries helps distinguish the direction of association from the scale of the price movements.

These correlations should be read within the observed period. They combine recovery, supply conditions and policy responses, and do not hold those forces constant. A correlation matrix helps describe the data, but cannot establish which mechanism generated the association. The regression comparison below checks the same question in inflation percentage-point units.

\clearpage
\subsection{Regression estimates and a worked calculation}

The OLS calculation uses centred cross-products. Across 27 quarters, unemployment sums to 239.7, so $\bar U=239.7/27=8.8778$. Its centred sum of squares is 231.5667. For annual headline inflation, the corresponding cross-product is 21.8443:
\begin{equation}
\hat\beta_H=\frac{\sum_q(U_q-\bar U)(\pi_{H,q}-\bar\pi_H)}{\sum_q(U_q-\bar U)^2}
=\frac{21.8443}{231.5667}=0.0943.
\end{equation}
The fitted annual lines are $\hat\pi_H=4.3747+0.0943U$ and $\hat\pi_C=3.7507+0.1228U$. These are associations within the sample; the intercept at zero unemployment has no practical interpretation here.

\begin{table}[H]
\centering\small
\caption{OLS estimates: the same unemployment observations in each comparison}
\label{tab:regression}
\begin{tabular}{llrrrrrr}
\toprule
Inflation & Sample & $n$ & $\hat\beta_H$ & $\hat\beta_C$ & $\hat\delta$ & $R_H^2$ & $R_C^2$\\
\midrule
YoY & Full & 27 & 0.094 & 0.123 & 0.028 & 0.076 & 0.137\\
YoY & Exclusion & 24 & -0.060 & 0.420 & 0.481 & 0.006 & 0.289\\
q/q SAAR & Full & 27 & 0.059 & 0.222 & 0.162 & 0.007 & 0.288\\
q/q SAAR & Exclusion & 24 & -0.113 & 0.321 & 0.434 & 0.004 & 0.133\\
\bottomrule
\end{tabular}
\tabnote{$\delta=\beta_C-\beta_H$. Slopes give inflation percentage points associated with one percentage point of unemployment. Nominal HC3 ranges for $\delta$ appear in Figure~\ref{fig:delta}; no significance stars are used.}
\end{table}


\fig{04_phillips_scatter.pdf}{Matched inflation--unemployment observations and OLS fits. Diamonds mark excluded quarters; dashed fits use only the exclusion sample's observed unemployment range.}{fig:scatter}{Matched-quarter datasets; calculations.}

The full annual slope difference is $0.1228-0.0943=+0.0285$. The q/q difference is also positive, at $+0.162$. Under the exclusion, headline slopes become slightly negative but core slopes remain positive. None of the four primary comparisons supports a more negative core relationship. Annual $R^2$ is 0.076 for headline and 0.137 for core: the improved fit concerns a positive slope. In a one-regressor model with an intercept, $R^2=r^2$, so $0.276^2\approx0.076$ checks the calculation.

\clearpage
\subsection{Uncertainty, influence and robustness}

Figure~\ref{fig:delta} displays the direct comparison with nominal 95\% HC3 ranges. Most ranges include zero. The annual exclusion range is positive, but simulation checks show that persistence and overlap can undermine nominal coverage. We therefore use these ranges to show uncertainty, without claiming a stable positive structural Phillips curve.

\fig{05_slope_difference.pdf}{Core-minus-headline slopes and nominal HC3 ranges. Negative values favour a more negative core slope. All four estimates are positive.}{fig:delta}{Direct difference regressions; calculations.}

\fig{06_influence.pdf}{Influence checks: centred unemployment variation (A) and 27 full-sample leave-one-quarter-out fits per outcome (B). Every core slope remains positive.}{fig:influence}{Unemployment influence and leave-one-quarter-out calculations.}

\begin{table}[H]
\centering\small
\caption{Robustness of unemployment slopes to sample and price-basket choices}
\label{tab:robustness}
\begin{tabularx}{\linewidth}{Xrrrr}
\toprule
Check & H: YoY & C: YoY & H: q/q & C: q/q\\
\midrule
Full sample & 0.094 & 0.123 & 0.059 & 0.222\\
Exclude three pandemic quarters & -0.060 & 0.420 & -0.113 & 0.321\\
Exclude all of 2020--2021 & -0.318 & 0.488 & -0.325 & 0.267\\
Full sample, calendar trend & 0.149 & 0.065 & 0.099 & 0.219\\
Exclude food only & 0.094 & 0.169 & 0.059 & 0.220\\
Core without housing & 0.094 & 0.191 & 0.059 & 0.448\\
\bottomrule
\end{tabularx}
\tabnote{H denotes headline; C denotes operational core or the stated alternative. All paired comparisons use identical observations and controls. Food-only exclusion retains fuel/light.}
\end{table}


The 2020Q2 observation accounts for 61.38\% of centred unemployment variation; the three excluded quarters account for 75.43\%. Its OLS leverage is 0.6509, calculated as $1/27+0.613816$, compared with average leverage $2/27=0.0741$. This measures how strongly its unemployment value can influence the fitted line; it does not establish that the observation is erroneous. Removing them reduces their influence but also sharply reduces the information available to estimate a slope. Food-only exclusion and removing housing leave the core comparison positive. With a calendar trend, the annual slope difference turns negative, but both slopes remain positive. The trend-adjusted result therefore does not provide evidence of a negative core relationship.

\clearpage
\subsection{How much headline variability is allocated to food?}

Food receives 88.71\% of monthly Combined headline-inflation variance and 80.96\% of annual variance. Urban allocations are lower, at 79.85\% and 64.69\%. All exceed food's respective basket weights of 45.86\% and 36.29\%. Food's variance allocation therefore exceeds its expenditure weight in both baskets and at both horizons.

\begin{table}[H]
\centering\small
\caption{Complete allocation of official headline-inflation variance (\%)}
\label{tab:variance}
\begin{tabular}{llrrrr}
\toprule
Basket & Change & Food & Fuel/light & Other four groups & Residual\\
\midrule
Combined & Monthly & 88.71 & 1.45 & 8.96 & 0.88\\
Combined & YoY & 80.96 & 4.03 & 14.94 & 0.07\\
Urban & Monthly & 79.85 & 1.57 & 18.98 & -0.41\\
Urban & YoY & 64.69 & 6.83 & 28.56 & -0.08\\
\bottomrule
\end{tabular}
\tabnote{Each row sums to 100\% before rounding. Monthly changes: February 2013--December 2025, $n=155$; YoY: January 2014--December 2025, $n=144$. These are covariance allocations, not causal effects.}
\end{table}


\fig{07_food_variance.pdf}{Food allocations to official raw headline variance; descriptive 12-month moving-block ranges use 4,999 draws.}{fig:variance}{Official CPI group contributions and paired moving-block calculations.}

For the monthly Combined basket, headline variance is 0.44393 and its covariance with the food contribution is 0.39381, both in squared percentage points. Applying equation~\eqref{eq:share},
\begin{equation}
s_F=100\left(\frac{0.39381}{0.44393}\right)=88.71\%.
\end{equation}
This calculation includes food's covariance with other contributions. It is not the variance of food inflation divided by headline variance. The residual closes differences between published headline and weighted group indices, so the complete allocation remains exact.

A component with a steady positive contribution can raise average headline inflation while explaining little of its variability. Conversely, a contribution that swings above and below its mean can receive a large variance share. The covariance formula captures those swings and their alignment with the total. For this reason, our shares should not be read as a breakdown of the average inflation rate, or as the percentage contribution to an individual month's price increase.

The denominator measures how much headline inflation varies around its own sample mean. Keeping the same denominator for every group is essential; using each group's variance instead would answer a different question.

The smaller Urban food weight helps explain the lower Urban allocation, but it is not the only reason: volatility and covariance also differ between baskets. The remaining categories and residual collectively account for 11.29\% of monthly Combined variance. This does not mean their prices or household welfare effects are unimportant. The result describes changes around mean inflation, not which component contributes most to every month's inflation rate.

The annual Combined descriptive range, 70.48--98.98\%, is wider than the monthly range, 85.94--93.14\%. Food remains dominant in these resampling checks, but the exact percentage should not be presented as a universal constant. Different horizons answer different questions.

\clearpage
\subsection{Monthly contributions and validation of the allocation}

The monthly contribution paths are consistent with the variance results. Food repeatedly tracks the larger swings in official headline inflation, while the combined non-food contribution is usually smaller (Figure~\ref{fig:contributions}). Food need not be the largest contribution in every month for its full-sample variance share to be high.

\fig{08_monthly_contributions.pdf}{Official monthly headline inflation and food contributions. Shading combines all other contributions and the residual; raw changes retain seasonality.}{fig:contributions}{Monthly exact-contribution datasets; calculations.}

\begin{table}[H]
\centering\small
\caption{Combined CPI: sensitivity of the monthly food allocation}
\label{tab:variancechecks}
\begin{tabularx}{\linewidth}{Xrr}
\toprule
Sample & $n$ & Food share (\%)\\
\midrule
Full history & 155 & 88.71\\
April 2018--December 2024 & 81 & 89.07\\
Exclude imputation-exposed changes & 152 & 88.61\\
February 2013--December 2019 & 83 & 87.11\\
January 2022--December 2025 & 48 & 90.27\\
\bottomrule
\end{tabularx}
\tabnote{Imputation exposure removes changes using April or May 2020 in either the numerator or denominator: April--June 2020. All calculations use the official raw headline target.}
\end{table}


The 89.07\% matched-period allocation is close to the full-history 88.71\%. Excluding imputation-exposed monthly changes gives 88.61\%, and the pre-2020 and post-2021 allocations remain high. These comparisons suggest that food dominance is not generated solely by a longer price sample or by the exceptional collection period.

Source validation supports this interpretation. Urban CPI contains 156 consecutive months and 1,092 observations across the groups and General index, with no missing levels. Fifty-six independent anchor comparisons cover eight published months. Weighted group levels reproduce official Urban headline within 0.08034 index points. Exact accounting retains a residual instead of forcing that gap to disappear.

Coverage, source agreement and accounting reconciliation are checked separately. None is used as a substitute for the other two.

An earlier 90.59\% estimate uses a reconstructed seasonally adjusted Combined aggregate. It answers a different accounting question. Independently adjusting General and each group with STL creates a large seasonal-reconciliation residual; Appendix~\ref{app:validation} records the distinction. Our main answer consistently uses official raw headline changes.


\clearpage
\subsection{Further calculations from monthly contributions}

Food's contribution has a slightly lower mean but about four times the standard deviation of the other contributions. These are percentage-point contributions to headline inflation, not inflation rates within the separate baskets. The other series includes the residual so that the two series add to official headline inflation in every month.

\begin{table}[H]
\centering\small
\caption{Monthly Combined CPI contributions: summary statistics}
\label{tab:contributionstats}
\begin{tabular}{lrrrrr}\toprule
Series & Mean & Median & SD & Minimum & Maximum\\\midrule
Food contribution & 0.199 & 0.218 & 0.605 & -1.656 & 2.635\\
All other contributions & 0.216 & 0.195 & 0.151 & -0.266 & 0.958\\
\bottomrule\end{tabular}
\tabnote{Percentage points; 155 raw monthly changes. Other contributions include the aggregation residual.}
\end{table}


\fig{11_calendar_contributions.pdf}{Combined CPI: calendar-month mean contributions. Each mean uses 12 or 13 observations from February 2013--December 2025.}{fig:calendar}{Calculations from the raw monthly contribution series.}

The mean food contribution is 0.941 percentage points in July and $-0.527$ in December. Its calendar pattern is larger than that of the other contributions. This describes the raw sample and does not identify weather or supply effects. Keeping seasonality in the main variance calculation is appropriate because the question concerns variation in the official unadjusted headline rate.

\fig{12_covariance_calculation.pdf}{Food allocation split into its own-variance and covariance terms. Both use official headline variance as the denominator.}{fig:covsplit}{Calculations from the same 155 monthly observations.}

Food's contribution variance is 0.366607 and its covariance with the other contributions is 0.027200. Their sum, 0.393807, is the covariance with headline inflation. Dividing each term by headline variance gives $82.58+6.13=88.71\%$. The positive covariance term shows that movements in the other contributions reinforce food's variance allocation on average. It is an accounting relationship, not evidence that one group causes price changes in another.

\clearpage
\subsection{Extension: food--non-food predictability}

The monthly extension uses 147 continuous observations, October 2013--December 2025. Non-food or core monthly log inflation is regressed on current food inflation, three food lags, one own lag, month effects and a pooled pandemic control. The delayed coefficient sum describes the combined lag association; a joint test asks whether all three delayed coefficients can be zero together. Opposite lag coefficients can offset in the sum.

\begin{table}[H]
\centering\small
\caption{Monthly extension: delayed food-inflation association}
\label{tab:spillover}
\begin{tabular}{lrrr}
\toprule
Outcome & Sum of lags 1--3 & Nominal HC3 range & Joint bootstrap tail\\
\midrule
Non-food & -0.0035 & $[-0.0592,\ 0.0523]$ & 0.4361\\
Core & -0.0188 & $[-0.0690,\ 0.0313]$ & 0.4661\\
\bottomrule
\end{tabular}
\tabnote{$n=147$, October 2013--December 2025. The joint test concerns all three delayed coefficients, not their sum alone. Restricted-null block-wild calculation: 19,999 draws, six-month blocks.}
\end{table}


\fig{09_spillover.pdf}{Delayed and current-plus-delayed food coefficients with nominal HC3 ranges. Blue excludes the current month; rust includes it.}{fig:spillover}{Monthly-model contrast estimates; calculations.}

The delayed sums, $-0.0035$ and $-0.0188$, do not detect a positive delayed association under this specification. Their ranges include zero and small positive values. Adding the current-month coefficient gives positive sums of 0.0342 and 0.0180, but their ranges also include zero. Common costs or simultaneous adjustment could generate a same-month association without delayed transmission.

Non-rejection does not establish zero spillover. In recursive simulations, a delayed sum of 0.048 is detected in only 5.4\% of the specified persistent-error non-food datasets. This is a scale comparison, not a replication of the RBI's 0.048 short-run ARDL sum~\cite{patra2024}. Reverse predictability is specification-dependent. These results show large direct food contributions to headline variability, with limited evidence of delayed predictability in the other price groups.

\section{Conclusion and policy implications}

The two assigned questions lead to different findings. Food receives about 89\% of monthly and 81\% of annual Combined headline-inflation variance. Urban allocations are about 80\% and 65\%. Yet excluding food and fuel does not uncover a stronger negative unemployment relationship in our matched Urban sample. Core is smoother at the quarterly horizon, but its unemployment slope is positive under both primary transformations and the fixed pandemic exclusion.

The results support examining headline and core together. Headline reflects the basket households face; core helps describe price pressures outside selected groups. A useful core measure need not reveal a particular Phillips-curve pattern. Equally, food's large allocation does not establish excessive demand or measure the effectiveness of a policy intervention.

Food-price monitoring should examine weather, arrivals, storage and distribution, consistent with RBI research~\cite{mohan2026,patra2024}. Our estimates do not identify an optimal policy response. Pandemic leverage, overlapping inflation and omitted expectations limit the unemployment comparison; low power limits the spillover findings. State-level and common-driver data are useful next steps.
\label{main:last}
% MAIN_TEXT_END

\par\medskip
\begin{minipage}{\linewidth}
\begingroup\setstretch{1.0}
\begin{thebibliography}{99}\fontsize{9}{10.8}\selectfont
\setlength{\itemsep}{2pt}\setlength{\parskip}{0pt}
\bibitem{cpiapi} MoSPI. \href{https://esankhyiki.mospi.gov.in/macroindicators?product=cpi&tab=api}{Consumer Price Index API, v2.0.0: historical 2012-base Urban series}. Accessed 21 November 2026.
\bibitem{cpirelease} NSO (2025). \href{https://www.mospi.gov.in/sites/default/files/press_release/CPI_PR_12Feb25.pdf}{Consumer Price Index Numbers, January 2025}. 12 February, Annexure I.
\bibitem{plfs2024} NSO. \href{https://mospi.gov.in/sites/default/files/publication_reports/PLFS_Quarterly_Bulletin_Oct-Dec24.pdf}{PLFS Quarterly Bulletin, October--December 2024}. Statement 5, Urban persons aged 15+.
\bibitem{plfs2019a} NSO. \href{https://www.mospi.gov.in/sites/default/files/publication_reports/Quarterly_Bulletin_PLFS_April_June_2019_M_0.pdf}{PLFS Quarterly Bulletin, April--June 2019}. Statement 5, p.~6 (PDF p.~9).
\bibitem{plfs2019b} NSO. \href{https://mospi.gov.in/sites/default/files/publication_reports/Quarterly_Bulletin_July_September_2019_M.pdf}{PLFS Quarterly Bulletin, July--September 2019}. Statement 5, p.~6 (PDF p.~9).
\bibitem{plfschange} MoSPI (2025). \href{https://www.mospi.gov.in/sites/default/files/publication_reports/PLFS_Changes-in-2025_rev.pdf}{PLFS: Changes in 2025}. Survey-design and comparability note.
\bibitem{imputation} NSO (2020). \href{https://www.mospi.gov.in/sites/default/files/press_release/CPI%20Technical%20Note%20on%20Imputation.pdf}{Compilation of Consumer Price Index: A Technical Note}. Table 1.
\bibitem{behera2017} Behera, H., Wahi, G. and Kapur, M. (2017). \href{https://rbi.org.in/Scripts/PublicationsView.aspx?id=17752}{Phillips Curve Relationship in India: Evidence from State-Level Analysis}. RBI Working Paper 08/2017.
\bibitem{patra2024} Patra, M.~D., John, J. and George, A.~T. (2024). \href{https://rbi.org.in/Scripts/BS_ViewBulletin.aspx?Id=22323}{Are Food Prices the `True' Core of India's Inflation?} RBI Bulletin, 18 January.
\bibitem{mohan2026} Mohan, R., Hasan, S., Sarkar, S. and John, J. (2026). \href{https://www.rbi.org.in/Scripts/BS_ViewBulletin.aspx?Id=23980}{Retail Inflation Volatility in India: Sources, Determinants, and Implications}. RBI Bulletin, 20 February.
\end{thebibliography}
\endgroup
\end{minipage}

\par\medskip
\appendix
\section{Regression details and inference checks}
\begin{table}[H]
\centering\small
\caption{Individual regression coefficients and nominal HC3 uncertainty}
\label{tab:regdetail}
\begin{tabular}{lllrrrr}
\toprule
Inflation & Sample & Outcome & Intercept & Slope & HC3 SE & Slope range\\
\midrule
YoY & Full & Headline & 4.375 & 0.094 & 0.039 & $[0.013,\ 0.175]$\\
YoY & Full & Core & 3.751 & 0.123 & 0.138 & $[-0.162,\ 0.408]$\\
YoY & Exclusion & Headline & 5.575 & -0.060 & 0.175 & $[-0.423,\ 0.302]$\\
YoY & Exclusion & Core & 1.400 & 0.420 & 0.126 & $[0.158,\ 0.682]$\\
q/q & Full & Headline & 4.676 & 0.059 & 0.115 & $[-0.178,\ 0.297]$\\
q/q & Full & Core & 2.695 & 0.222 & 0.040 & $[0.139,\ 0.305]$\\
q/q & Exclusion & Headline & 6.042 & -0.113 & 0.361 & $[-0.863,\ 0.636]$\\
q/q & Exclusion & Core & 1.938 & 0.321 & 0.139 & $[0.033,\ 0.609]$\\
\bottomrule
\end{tabular}
\tabnote{Nominal 95\% t-based ranges. They are descriptive sensitivity summaries in this short, dependent sample. Full and exclusion sample sizes are 27 and 24.}
\end{table}

\begin{table}[H]
\centering\small
\caption{Null rejection rates in recursive monthly-model simulations (\%)}
\label{tab:calibration}
\begin{tabularx}{\linewidth}{lXrrr}
\toprule
Outcome & Error process & HAC(6) & HC3 & Block wild\\
\midrule
Non-food & Independent Gaussian & 13.3 & 5.1 & 4.0\\
Non-food & AR/GARCH Gaussian & 14.7 & 5.4 & 5.1\\
Non-food & AR/GARCH, heavy tails & 13.7 & 5.9 & 6.1\\
Core & Independent Gaussian & 13.4 & 5.1 & 4.3\\
Core & AR/GARCH Gaussian & 15.9 & 5.9 & 7.1\\
Core & AR/GARCH, heavy tails & 14.2 & 5.4 & 5.9\\
\bottomrule
\end{tabularx}
\tabnote{Nominal rate: 5\%. Each cell uses 1,000 datasets; bootstrap tests use 999 inner draws. The Monte Carlo SE near 5\% is about 0.7 percentage points.}
\end{table}

\fig{10_power.pdf}{Power for the non-food equation under AR/GARCH errors. Each point uses 1,000 recursive datasets and 999 bootstrap draws. Rejection is 5.4\% at a delayed sum of 0.048, 31.6\% at 0.100 and 93.8\% at 0.200. Higher null rejection can inflate apparent power.}{fig:power}{Recursive simulation; true effects allocated equally over three food lags.}

\Needspace{6\baselineskip}
\section{Accounting and source checks}
\label{app:validation}

The earlier 90.5899\% share uses a reconstructed aggregate of separately seasonally adjusted groups. It differs from the raw official-headline target. Adjusting General and every group independently with STL assigns 24.2410\% of Combined variance and 25.0889\% of Urban variance to the reconciliation residual. Corresponding food shares are 69.7544\% and 61.3913\%. This residual includes seasonal-adjustment inconsistencies and is not an economic category.

PLFS corrections use Statement 5, printed p.~6 (PDF p.~9), in the two 2019 bulletins. CPI imputation flags follow the technical note: Urban pan/tobacco is macro-imputed in April 2020 but is not marked in May; Urban housing and fuel/light are not marked in either month. Elementary-level imputation may still occur. Annual exclusions remove changes using April--May 2020 in either period, including April--May 2021.

Calendar HAC respects actual quarter distances after exclusions. The monthly simulations regenerate lagged outcomes recursively. Block-wild null rejection ranges from 4.0--6.1\% for non-food and 4.3--7.1\% for core across the simulated error processes. These checks do not guarantee correct inference under every process. Reverse predictability remains exploratory: adding a lagged relative-price term changes one joint bootstrap tail from 0.0965 to 0.0341, demonstrating specification sensitivity.
\end{document}
